\documentclass[11pt,a4paper]{article}
\usepackage[margin=2.4cm]{geometry}
\usepackage{amsmath,amssymb,bm,mathtools}
\usepackage{enumitem}
\usepackage[colorlinks=true,linkcolor=blue,citecolor=blue]{hyperref}

\newcommand{\bw}{\bm{\omega}}
\newcommand{\bq}{\mathbf{q}}
\newcommand{\bj}{\mathbf{j}}
\newcommand{\bh}{\mathbf{h}}
\newcommand{\bk}{\mathbf{k}}
\newcommand{\bN}{\mathbf{N}}
\newcommand{\bGamma}{\bm{\Gamma}}
\newcommand{\It}{\tilde{\mathbf{I}}}
\newcommand{\E}{\mathsf{E}}
\newcommand{\Id}{\mathsf{1}}
\newcommand{\K}{\mathsf{K}}
\newcommand{\Lm}{\mathsf{L}}
\newcommand{\se}{\sigma_{\mathrm e}}
\newcommand{\Aw}{A_{\mathrm w}}
\newcommand{\dd}{\,\mathrm{d}}

\title{The gyroscopic matrix $\bN$:\\ how the fluid core and damping enter}
\author{Working notes (prepared with Claude) --- to be checked}
\date{1 October 2026}

\begin{document}
\maketitle

\noindent This note expands on the matrix $\bN$ of \S10.1 of \texttt{rotation\_dynamic\_elastic.pdf}
(equation numbers prefixed by ``N'' refer to that note). $\bN$ multiplies $\dot\bw$ in the rotational
row of the weak form,
\begin{equation}
  (\It\bw+\bN\dot\bw-\bj+\bk)\cdot\bw'=0,
  \label{eq:row}
\end{equation}
and carries everything about the inertia of the rotation vector itself. The fluid core and the
damping each change $\bN$, in different ways: the core rescales it, and damping adds a new part.
Throughout, $\E=\begin{psmallmatrix}0&-1\\1&0\end{psmallmatrix}$, with $\E^2=-\Id$, $\E^T=-\E$ and
$\E^{-1}=-\E$. The equatorial (top-left $2\times2$) block of $\bN$ is written $\bN_\perp$.

%=======================================================================
\section{Where $\bN$ comes from: no core, no damping}
%=======================================================================

\begin{enumerate}
  \item Start from the equatorial Liouville equation, eq.~N(10), with $\bh=\bGamma=\mathbf{0}$:
  \begin{equation}
    I_1\dot\bw+\dot\bj=\Omega\E\big[(I_3-I_1)\bw-\bj\big].
  \end{equation}
  \item Drop $\dot\bj$, which is $O(\nu/\Omega)$ relative to the retained terms (N\S4.2):
  \begin{equation}
    I_1\dot\bw=\Omega\E\big[(I_3-I_1)\bw-\bj\big].
  \end{equation}
  \item Multiply both sides on the left by $\E^{-1}/\Omega=-\E/\Omega$. On the right this removes $\Omega\E$:
  \begin{equation}
    -\frac{I_1}{\Omega}\E\dot\bw=(I_3-I_1)\bw-\bj .
  \end{equation}
  \item Add $(I_1/\Omega)\E\dot\bw+\bj$ to both sides:
  \begin{equation}
    \boxed{\;(I_3-I_1)\bw+\frac{I_1}{\Omega}\E\,\dot\bw=\bj\;}
  \end{equation}
  \item Comparing with eq.~\eqref{eq:row} (with $\bk=\mathbf{0}$), the equatorial block of $\It$ is
  $(I_3-I_1)\Id$ and
  \begin{equation}
    \bN_\perp=\frac{I_1}{\Omega}\,\E .
  \end{equation}
  This is a multiple of $\E$, and therefore antisymmetric.
  \item \textbf{Axial part.} The third row and column of $\bN$ are zero, so the axial component
  satisfies $-I_3\omega_3=j_3$. This is algebraic: $\omega_3$ has no memory.
  \item \textbf{Static limit.} Setting $\dot\bw=\mathbf{0}$ recovers AA24's static relation $\It\bw=\bj$.
\end{enumerate}

%=======================================================================
\section{The fluid core: it rescales $\bN$}
%=======================================================================

\begin{enumerate}
  \item With the core, the relative angular momentum is no longer zero. To leading order in $\nu/\Omega$
  (N\S4.1),
  \begin{equation}
    \bh_{\rm core}\approx\frac{A_f}{\Omega}\,\E\dot\bw,
  \end{equation}
  where $A_f$ is the fluid core's equatorial moment. The core's rotation differs from the mantle's by
  only $O(\epsilon)$ of $\bw$.
  \item This is small, but it enters the Liouville equation, eq.~N(10), through the term $-\Omega\E\bh$:
  \begin{equation}
    I_1\dot\bw+\dot\bh=\Omega\E\big[(I_3-I_1)\bw-\bj\big]-\Omega\E\bh .
  \end{equation}
  \item Evaluate the new term. The factor $\Omega$ cancels the $1/\Omega$, and $\E^2=-\Id$ flips the sign:
  \begin{equation}
    -\Omega\E\bh_{\rm core}=-\Omega\E\,\frac{A_f}{\Omega}\E\dot\bw=-A_f\E^2\dot\bw=+A_f\dot\bw .
  \end{equation}
  \item Move it to the left-hand side:
  \begin{equation}
    (I_1-A_f)\dot\bw+\dot\bh_{\rm core}=\Omega\E\big[(I_3-I_1)\bw-\bj\big].
  \end{equation}
  \item $\dot\bh_{\rm core}=(A_f/\Omega)\E\ddot\bw$ is a further factor $\nu/\Omega$ smaller than
  $A_f\dot\bw$, so drop it. With $A_m:=I_1-A_f$, the mantle's equatorial moment,
  \begin{equation}
    A_m\dot\bw=\Omega\E\big[(I_3-I_1)\bw-\bj\big].
  \end{equation}
  \item Repeating steps 3--5 of \S1 with $I_1$ replaced by $A_m$:
  \begin{equation}
    \boxed{\;\bN_\perp=\frac{A_m}{\Omega}\,\E\;}
  \end{equation}
  \item \textbf{Interpretation.} The core only rescales $\bN$. It stays a multiple of $\E$, and so stays
  antisymmetric. Writing $\Aw$ for ``$I_1$ or $A_m$'', and allowing a triaxial Earth, the full matrix is
  \begin{equation}
    \bN=\frac{1}{\Omega}\begin{pmatrix}0&-A_{w,2}&0\\A_{w,1}&0&0\\0&0&0\end{pmatrix}.
  \end{equation}
  \item Nothing else in the problem changes. The static limit is untouched, since it does not involve
  $\bN$, and the Euler frequency becomes $\se=\Omega(I_3-I_1)/\Aw$. PREM gives $A_f/I_1=0.113$.
\end{enumerate}

%=======================================================================
\section{Damping: it adds a symmetric part to $\bN$}
%=======================================================================

\begin{enumerate}
  \item The solver uses the damped law (N\S8.1)
  \begin{equation}
    \dot\bw=\Lm(\bw-\bq),\qquad \Lm=\se\E-\gamma\Id,\qquad \bq=\frac{\bj}{I_3-I_1}.
  \end{equation}
  To put it in the form of eq.~\eqref{eq:row}, invert $\Lm$.
  \item Multiply $\Lm$ by $-\se\E-\gamma\Id$, using $\E^2=-\Id$:
  \begin{equation}
    (\se\E-\gamma\Id)(-\se\E-\gamma\Id)=-\se^2\E^2-\se\gamma\E+\se\gamma\E+\gamma^2\Id=(\se^2+\gamma^2)\Id .
  \end{equation}
  \item Hence
  \begin{equation}
    \Lm^{-1}=-\frac{\se\E+\gamma\Id}{\se^2+\gamma^2}.
  \end{equation}
  \item Multiply the damped law on the left by $\Lm^{-1}$: $\Lm^{-1}\dot\bw=\bw-\bq$.
  \item Multiply by $-(I_3-I_1)$ and use $(I_3-I_1)\bq=\bj$:
  \begin{equation}
    \boxed{\;(I_3-I_1)\bw+\frac{(I_3-I_1)(\se\E+\gamma\Id)}{\se^2+\gamma^2}\,\dot\bw=\bj\;}
  \end{equation}
  So with damping,
  \begin{equation}
    \bN_\perp=\frac{(I_3-I_1)(\se\E+\gamma\Id)}{\se^2+\gamma^2}.
  \end{equation}
  \item \textbf{Check.} At $\gamma=0$: $\bN_\perp=(I_3-I_1)\E/\se=(I_3-I_1)\E\,\Aw/[\Omega(I_3-I_1)]=(\Aw/\Omega)\E$,
  as in \S\S1--2.
  \item \textbf{Split into two parts.} $\bN_\perp=\bN_a+\bN_s$, with
  \begin{align}
    \bN_a&=\frac{(I_3-I_1)\se}{\se^2+\gamma^2}\,\E=\frac{\Aw}{\Omega}\,\frac{1}{1+\gamma^2/\se^2}\,\E
    &&\text{(antisymmetric, gyroscopic)},\\
    \bN_s&=\frac{(I_3-I_1)\gamma}{\se^2+\gamma^2}\,\Id
    &&\text{(symmetric, dissipative)}.
  \end{align}
  \item \textbf{Sizes.} With $\gamma=\se/(2Q)$, $\gamma/\se=1/(2Q)=0.005$ for $Q=100$.
  \begin{itemize}
    \item $\bN_a$ is the undamped $\bN$ reduced by the factor $1/(1+\gamma^2/\se^2)$, i.e.\ by $2.5\times10^{-5}$.
    This is negligible.
    \item $\bN_s$ is new. Its size relative to $\bN_a$ is $\gamma/\se=1/(2Q)$.
  \end{itemize}
\end{enumerate}

%=======================================================================
\section{Why the split matters: an energy identity}
%=======================================================================

\begin{enumerate}
  \item Consider free motion: no direct load, so $\bj=(I_3-I_1)\K\bw$ (N\S6.2). The equatorial
  rotational row becomes
  \begin{equation}
    (I_3-I_1)(\Id-\K)\bw+\bN\dot\bw=\mathbf{0}.
  \end{equation}
  \item Take the dot product with $\dot\bw$:
  \begin{equation}
    (I_3-I_1)\,\dot\bw^T(\Id-\K)\bw+\dot\bw^T\bN_a\dot\bw+\dot\bw^T\bN_s\dot\bw=0.
  \end{equation}
  \item The antisymmetric part does no work. $\dot\bw^T\bN_a\dot\bw$ is a number, so it equals its
  transpose, $\dot\bw^T\bN_a^T\dot\bw=-\dot\bw^T\bN_a\dot\bw$. It is therefore zero.
  \item $\Id-\K$ is symmetric, so $\dot\bw^T(\Id-\K)\bw=\frac{\dd}{\dd t}\big[\tfrac12\bw^T(\Id-\K)\bw\big]$.
  \item Hence
  \begin{equation}
    \boxed{\;\frac{\dd}{\dd t}\Big[\tfrac12(I_3-I_1)\,\bw^T(\Id-\K)\bw\Big]=-\dot\bw^T\bN_s\dot\bw\le0\;}
  \end{equation}
  The inequality holds because $\bN_s$ is a positive multiple of $\Id$ when $\gamma>0$.
  \item \textbf{Consequences.}
  \begin{itemize}
    \item With $\gamma=0$ the bracketed ``energy'' is conserved. This is the conserved quantity that
    makes the free orbit an ellipse in the detailed derivation note
    of \S7.
    \item With $\gamma>0$, $\bN_s$ drains it, so the wobble decays.
    \item The core changes only $\bN_a$. It therefore changes the wobble frequency but not the decay.
  \end{itemize}
\end{enumerate}

%=======================================================================
\section{Why the main note writes things in terms of $\bN$}
%=======================================================================

\begin{enumerate}
  \item The reciprocity argument (N\S10.2) needs only $\bN$. The adjoint problem uses $\bN^\dagger=-\bN^T$.
  \item Transposing flips the sign of the antisymmetric part and keeps the symmetric part, so
  \begin{equation}
    -\bN^T=-(\bN_a+\bN_s)^T=\bN_a-\bN_s .
  \end{equation}
  \item Run backwards in time, $\dd/\dd t\to-\dd/\dd\tau$, this becomes $-\bN_a+\bN_s$ (N\S10.3):
  \begin{itemize}
    \item the gyroscopic part is reversed, i.e.\ the rotation is reversed;
    \item the dissipative part is kept, so the adjoint is damped in reversed time just as the
    forward problem is in forward time.
  \end{itemize}
  \item Both the core (a rescaled $\bN_a$) and the damping (an added $\bN_s$) therefore fit the same
  reciprocity theorem, with no new derivation needed.
  \item In the forward solver itself, the equivalent form $\dot\bw=\Lm(\bw-\bq)$, with $\se$ built from
  $\Aw$, is simpler to use (N\S8). The two forms describe the same equation.
\end{enumerate}

\end{document}
